\SourceLeaf{44}expanded, the general term is $r^{\text{tum}} = \cos\left(2r+\frac{1-(-1)^n}{2}\right)\varphi$. For when $n$ is even this form clearly reduces to $\cos 2r\varphi$, and when it is odd it becomes $\cos(2r+1)\varphi$.

In the short treatise itself I used this general notation only in the final formula (\S\,11). It could be profitably applied already in the investigations leading to that formula. I thought I should refrain from using it there, since the additional remarks required on the transformations of such signs when they enter the calculations themselves would have led to undue length and taken me beyond the limits of the treatise. I thought it better to devote a separate investigation to the treatment of such signs when the terms they affect enter lengthy calculations, and to the devices needed to keep those calculations concise.

1. At first the integration of the proposed formula seems very easy to carry out. The given quantity $(1+n\cos\varphi)^\nu$ can, by the binomial theorem, be converted into a series proceeding by successive integral powers of $\cos\varphi$. Then its individual terms, multiplied by $d\varphi$ and integrated, will give the desired result. This is indeed correct, but proceeding in this way leads us to very complicated expressions. Each part of the integral sought will be represented by a rather lengthy series. It is therefore preferable to follow another method, based on the fact that, by the rules of analytical trigonometry, any power of $\cos\varphi$ with a positive integral exponent can be expanded into a very concise series whose terms contain only cosines of multiples of $\varphi$. The proposed task will therefore be carried out by first expressing $(1+n\cos\varphi)^\nu$ as a series proceeding by powers of $\cos\varphi$, then transforming each term of that series, according to the general law established beforehand, into a new series proceeding by cosines of multiples. Finally, from all these series we collect the terms containing the cosine of the same angle. The resulting new series, duly ordered and multiplied by $d\varphi$, can be integrated at once.

2. First, therefore, we must examine the law by which a power of a given cosine is expressed in cosines of multiples of the angle. This investigation is based on the elementary trigonometric formula\HistoricalNote{*)}{*)}{\textsc{Kaestner}, Astron. Abh., Erste Sammlung, I, 10.}, which states that
